\section{局域场因子、STLS 与有效相互作用}

超越 RPA 的系统化方式是局域场因子（GV §5.4--5.6）。

\subsection{定义}

\begin{equation}
  \chi_{\mathrm{irr}}
  =\frac{\chi_0}{1+v_{\mathbf{q}}G(\mathbf{q},\omega)\chi_0},
  \qquad
  \varepsilon
  =1-\frac{v_{\mathbf{q}}\chi_0}{1+v_{\mathbf{q}}G\chi_0}.
\label{eq:local_field}
\end{equation}
自旋通道可定义 $G_\pm$ 或 $G_{\uparrow\uparrow},G_{\uparrow\downarrow}$。压缩率与自旋磁化率增强
\begin{equation}
  \frac{K}{K_0}
  =\frac{1}{1-v_{\mathrm{scr}}N(0)(1-G_+(0))},
\end{equation}
（示意；精确预因子见 GV）把 $G(q\to0)$ 与热力学相连：
\begin{equation}
  \lim_{q\to0}G_+(q,0)
  =1-\frac{K_0}{K}.
\end{equation}

\subsection{Hubbard 与 STLS}

Hubbard 近似用交换孔的特征波矢插值
\begin{equation}
  G_H(q)
  \approx
  \frac{q^2}{q^2+k_F^2}
  \quad\text{（三维示意）},
\end{equation}
使 $q\to\infty$ 与 $q\to0$ 极限合理。STLS（Singwi--Tosi--Land--Sj\"olander）用
\begin{equation}
  G(q)
  =-\frac{1}{n}
  \int\frac{\mathrm{d}^d q'}{(2\pi)^d}
  \frac{\mathbf{q}\cdot\mathbf{q}'}{q'^2}
  \bigl[S(\mathbf{q}-\mathbf{q}')-1\bigr]
\end{equation}
与 $S(q)$ 通过涨落耗散自洽，显著改善 $g(0)$ 与关联能。Slater 型 $G_{\uparrow\uparrow}$ 由 $S_0-1$ 给出，$G_{\uparrow\downarrow}=0$（纯交换）。

\subsection{有效相互作用}

三种常用屏蔽势（GV §5.5）：
\begin{itemize}
  \item 试探电荷--试探电荷：$W_{tt}=v_q/\varepsilon$；
  \item 电子--试探电荷：多一个 $(1-G)$ 因子；
  \item 电子--电子：含 $G_\pm$，进入准粒子散射与超导配对核。
\end{itemize}

\subsection{Friedel 振荡}

杂质静势的诱导密度
\begin{equation}
  \delta n(r)
  \propto
  \frac{\cos(2k_F r+\eta)}{r^{d}}
\end{equation}
（渐近）。振荡波矢由 $\chi(q,0)$ 在 $2k_F$ 的奇异决定；振幅依赖 $G(q)$。STM 可探测局域态密度的相应振荡。

\begin{example}[对照 GV]
静态 $F(q)$、$G(q/k_F)$、等离色散与 $S(q,\omega)$ 的图是电子液理论的“标准像”。建议用 \eqref{eq:Fstd}、\eqref{eq:RPA}、\eqref{eq:ERPA} 复现高密极限，再讨论 $g(0)$ 问题。
\end{example}
